% gz-p10-04-1: 极值点判别——y = x^3 - 3x，标极大 x=-1，极小 x=1，及 f' 符号
\documentclass[tikz,border=4pt]{standalone}
\usepackage{ctex}
\usepackage{amsmath}
\usepackage{pgfplots}
\pgfplotsset{compat=1.18}

\begin{document}
\begin{tikzpicture}
\begin{axis}[
  width=11cm, height=8cm,
  axis lines=center,
  xlabel={$x$}, ylabel={$y$},
  every axis x label/.style={at={(axis cs:2.4,0)}, anchor=west},
  every axis y label/.style={at={(axis cs:0,3.6)}, anchor=south},
  xmin=-2.4, xmax=2.4,
  ymin=-3.6, ymax=3.6,
  xtick={-1,1},
  ytick={-2,2},
  tick style={color=black},
  ticklabel style={font=\small},
  clip=true,
]
  % 曲线 y = x^3 - 3x
  \addplot[black, thick, domain=-2.2:2.2, samples=120] {x^3 - 3*x};
  \node[right, black] at (axis cs:1.8, -1.8) {$y=x^3-3x$};

  % 极大值点 (-1, 2)
  \addplot[only marks, mark=*, mark size=2.5pt, red] coordinates {(-1,2)};
  \node[above, red, font=\small] at (axis cs:-1,2) {极大值};
  \node[left, red, font=\footnotesize] at (axis cs:-1, 2) {$(-1,2)$};

  % 极小值点 (1, -2)
  \addplot[only marks, mark=*, mark size=2.5pt, blue] coordinates {(1,-2)};
  \node[below, blue, font=\small] at (axis cs:1,-2) {极小值};
  \node[right, blue, font=\footnotesize] at (axis cs:1,-2) {$(1,-2)$};

  % 虚线
  \draw[gray, dashed] (axis cs:-1,2) -- (axis cs:-1,0);
  \draw[gray, dashed] (axis cs:1,-2) -- (axis cs:1,0);

  % f' 符号标注
  \node[font=\footnotesize, red] at (axis cs:-1.8, 3.0) {左:$f'>0$};
  \node[font=\footnotesize, red] at (axis cs:-0.4, 3.0) {右:$f'<0$};
  \node[font=\footnotesize, blue] at (axis cs:0.4, -3.2) {左:$f'<0$};
  \node[font=\footnotesize, blue] at (axis cs:1.8, -3.2) {右:$f'>0$};
\end{axis}
\end{tikzpicture}
\end{document}
